\documentclass[uplatex,dvipdfmx]{jlreq}
\usepackage{jlreq-deluxe}
\usepackage[noalphabet]{pxchfon}
\usepackage{stix2}
\usepackage[fleqn,tbtags]{mathtools}
\usepackage{hira-stix}
\usepackage{float}
\usepackage{booktabs}
\usepackage{multirow}
\usepackage{siunitx}
\usepackage{listings}
\usepackage{xcolor}
\usepackage[hyphens]{url}
\usepackage{array}
\usepackage{longtable}

% ===== コードリスト設定（必要に応じて変更） =====
\lstset{
    basicstyle=\ttfamily\small,
    keywordstyle=\color{blue}\bfseries,
    commentstyle=\color{gray}\itshape,
    stringstyle=\color{orange},
    numbers=left,
    numberstyle=\tiny\color{gray},
    stepnumber=1,
    numbersep=10pt,
    frame=single,
    frameround=tttt,
    breaklines=true,
    breakatwhitespace=true,
    tabsize=2,
    captionpos=b,
    showstringspaces=false
}

\begin{document}

% ===== ヘッダー情報 =====
\title{\textbf{第6週目事前課題 作業ログ}}
\author{25G1144 脇阪隼人}
\date{\today}
\maketitle

\noindent
\textbf{Arduinoオーケストラ}

\bigskip

% ===== セクション1: 進捗サマリー =====
\section{進捗サマリー（要約）}

\begin{itemize}
    \item \textbf{全体進捗率：}Ptype2でSA／SPの和音対応・演奏順不具合修正・中級アレンジ譜への更新が完了し、Ptype3として16分音符単位のTick送信方式への移行（整合性修正）も完了。実機書き込み・通し演奏での最終検証が残っている。
    \item \textbf{マイルストーン達成状況：}達成（Slave Arduino・Slave Processingの主要機能追加分、Ptype3でのTick解像度変更分）
    \item \textbf{主要成果：}
    \begin{itemize}
        \item SA\_ptype1.ino のコンパイルエラー（配列リファクタの不整合、UNO R4 WiFiでの volatile String 起因エラー）を特定・修正
        \item SA／SPで主旋律と対旋律を同時演奏できる機能を追加（4値シリアル送信＋Minimでのミックス再生）
        \item SA に myPlayOrder==4（ドラム）の発音ロジックを追加。ハイハット／スネアを偶数Tickごとに交互発音
        \item DrumProcessing\_new.pde を新しいSA仕様に合わせて修正し、未インストールライブラリ（processing.sound）依存とAPI誤用（Noise.setAmplitude）を解消
        \item ファイル構成をMA/MP1・MP2/SA/SP\_Piano・SP\_Drumに再編し、GitHubへpush
        \item Ptype2のSA楽譜配列を「主旋律1音」から「主旋律・対旋律それぞれ最大3音の和音」対応の二次元配列に拡張し、送信フォーマットを4値（音階0,1,2,音の長さ）に変更
        \item MP1/MP2/else.MP\_ptype2のinstNames順序不一致（木琴とフルートの演奏順が入れ替わって送信される不具合）を特定・修正
        \item もりのくまさん中級アレンジ譜（gakuhu\_tyukyu.pdf）の全13小節を採譜し、SAの楽譜配列を\texttt{pitch\_to}/\texttt{pitch\_he}（104要素）に全面更新
        \item Ptype3を作成し、楽譜解像度を8分音符基準から16分音符基準（208要素）へ拡張した上で、MA/SA/SP\_DrumのTick送信間隔・演奏順の開始・合流タイミング・ドラム発音間隔・総発音数を整合させる修正を実施
    \end{itemize}
    \item \textbf{重大課題（影響度：高）：}
    \begin{itemize}
        \item ローカル環境にarduino-cli等のコンパイラがなく、Claude Codeでの修正はコードレビューのみで実機・実コンパイル未検証
        \item MA/MP側（Master Arduino／Processing）はメンバーが独自に大幅改修中で、SA/SPとのプロトコル整合性は未確認
        \item ピアノ実機ボードのINST\_ID設定が誤っている可能性（演奏順1番手で送信しているにもかかわらず輪奏部分を演奏する現象）を特定したが、実機側の\texttt{\#define INST\_ID}値の確認・修正はユーザー側作業として未対応
    \end{itemize}
\end{itemize}

% ===== セクション2: 作業ログ（詳細トラッキング） =====
\section{作業ログ（詳細トラッキング）}

% 各作業項目は \subsection で記述する
% 必要に応じて \subsubsection を追加してよい

\subsection{SA\_ptype1.ino のコンパイルエラー修正}
\noindent
作業日時：2026年6月18日 \quad 作業時間；4時間
状態：完了 \quad 進捗率：100\%

\noindent
\textbf{依存関係・ブロッカー：}なし

\noindent
\textbf{作業内容：}

楽譜配列を2次元（和音用）から1次元へリファクタリングした際、myPlayOrder==2/3の分岐が
削除済みの \texttt{pitch\_sub}/\texttt{duration\_sub}/\texttt{SCORE\_LENGTH\_SUB} を参照したままになっており、
かつ \texttt{pitch\_main[index][0]} のような2次元添字アクセスが残っていたためビルド不能になっていた。
単一の \texttt{pitch[]}/\texttt{duration[]} 配列と \texttt{startTick} オフセットによるロジックに統一して解消した。

その後、UNO R4 WiFi（Renesasコア）に書き込んだ際に
\texttt{passing 'volatile arduino::String' as 'this' argument discards qualifiers} が発生。
RenesasコアのStringクラスにはvolatile修飾のメンバ関数が無いことが原因であったため、
受信バッファを \texttt{volatile char} 配列＋長さ変数に置き換え、\texttt{loop()} 側で通常のStringへ
変換する方式に変更して解消した。

\subsection{演奏順（myPlayOrder）ごとの開始位置・合流ロジックの調整}
\noindent
作業日時：2026年6月18日 \quad 作業時間：2時間
状態：完了 \quad 進捗率：100\%

\noindent
\textbf{依存関係・ブロッカー：}なし

\noindent
\textbf{作業内容：}

myPlayOrder==2,3 が演奏を開始するTickと、楽譜上のどの要素目から始めるかを
\texttt{startTick}／\texttt{indexOffset} という2変数で一般化。要望に応じて開始要素を
37→33要素目に調整した後、TickCount=66（myPlayOrder==1が65要素目を演奏するタイミング）
以降はmyPlayOrder==1と同じindex進行に合流するよう変更した。

\subsection{主旋律・対旋律の同時演奏機能の追加（SA／SP）}
\noindent
作業日時：2026年6月18日 \quad 作業時間：2時間
状態：完了 \quad 進捗率：100\%

\noindent
\textbf{依存関係・ブロッカー：}対旋律データは過去にrevertされたコミット（対旋律の楽譜を作成）の
\texttt{pitch\_sub}/\texttt{duration\_sub} を再利用することで合意

\noindent
\textbf{作業内容：}

SA\_ptype1.ino で主旋律（\texttt{pitch\_main}/\texttt{duration\_main}）と対旋律
（\texttt{pitch\_sub}/\texttt{duration\_sub}）を毎Tickで計算し、
\texttt{mainPitch,mainDuration,subPitch,subDuration} の4値をSerialで送信するよう変更。
SP\_ptype1.pde（現SP\_Piano）側は \texttt{playVoice()} を追加し、2音をそれぞれ
\texttt{out.playNote()} に渡すことで、Minimの出力バス上でミックスして同時再生できるようにした。

\subsection{ドラム機能（myPlayOrder==4 / INST\_ID==4）の追加}
\noindent
作業日時：2026年6月18日 \quad 作業時間：1時間
状態：完了 \quad 進捗率：100\%

\noindent
\textbf{依存関係・ブロッカー：}既存の \texttt{Drumslave\_ver2.ino} を参考実装として利用

\noindent
\textbf{作業内容：}

SA\_ptype1.ino に \texttt{handleDrum()} を追加し、INST\_ID==4のときだけ通常楽器の処理から
分岐させ、他楽器（INST\_ID!=4）の既存ロジックには影響を与えないようにした。
要件：(1)ハイハット(drumType=2,vel=80)とスネア(drumType=1,vel=110)の2種類、
(2)\texttt{i2cReceiveCount-1} をTickCount（0始まり）として扱う、
(3)TickCountが偶数のときだけ発音、(4)発音ごとにハイハット→スネアと交互、
(5)TickCount==97（i2cReceiveCount==98）でisPlaying=falseにして終了、
(6)出力フォーマットは \texttt{Drumslave\_ver2.ino} と同じ \texttt{drumType,velocity,BPM}、
をすべて満たす実装とした。

\subsection{DrumProcessing\_new.pde のSA仕様対応・エラー修正}
\noindent
作業日時：2026-06-18 \quad 状態：完了 \quad 進捗率：100\%

\noindent
\textbf{依存関係・ブロッカー：}なし

\noindent
\textbf{作業内容：}

旧DrumProcessing\_new.pdeは\texttt{Drumslave\_ver2.ino}向けに書かれており、
\texttt{START}/\texttt{END}/\texttt{"DrumSlave ready"}という明示的なステータス文字列を前提にしていたが、
新しいSAの\texttt{handleDrum()}はこれらを送信せず、デバッグ用の行（\texttt{"I2C: ..."}等）も
同じシリアルに混在する。カンマ区切りでちょうど3項目の行のみを演奏データとして扱うように修正し、
\texttt{isPlaying}/\texttt{hitCount}の管理もステータス文字列に依存しない自己完結型のロジックに変更した
（偶数Tickのみ発音のため、全98Tick中の総発音数は49回であることも反映）。

その後、実行時に2つのエラーが発生：
\begin{itemize}
    \item \texttt{No library found for processing.sound}：Processing Soundライブラリが未インストール。
    本プロジェクトの他のSPファイルで既に動作確認済みの\texttt{ddf.minim}に置き換え、
    新規ライブラリのインストールを不要にして解消。
    \item \texttt{The function setAmplitude(float) does not exist}：Minimの\texttt{Noise}クラスには
    \texttt{setAmplitude()}メソッドが存在せず、振幅はコンストラクタで指定する仕様であったため、
    \texttt{new Noise(amplitude, Noise.Tint.WHITE)}の形に修正して解消。
\end{itemize}

\subsection{SA／SPの楽譜構造を和音対応の二次元配列化し、送信フォーマットを4値に拡張}
\noindent
作業日時：2026年6月23日 \quad 作業時間：4時間
状態：完了 \quad 進捗率：100\%

\noindent
\textbf{依存関係・ブロッカー：}対旋律(\texttt{pitch\_sub}/\texttt{duration\_sub})の実データは、過去にrevertされたコミット（対旋律の楽譜を作成）のものを復元して使用

\noindent
\textbf{作業内容：}

AIに楽譜画像をpdf形式で渡し、画像解析によって得た楽譜のmidi番号や持続時間を配列として作成してもらった。
実際には楽譜仕様通りの演奏を行うことができなかったため、独自に新しく2次元楽譜配列を作成した。
SAの楽譜配列を「主旋律1音のみ」の1次元配列から、主旋律\texttt{pitch\_main}・対旋律\texttt{pitch\_sub}それぞれが1要素につき最大3音を持つ二次元配列に拡張。
音の長さ配列（\texttt{duration\_main}/\texttt{duration\_sub}）は1次元のまま維持した。
myPlayOrder==1,2,3のときのシリアル送信フォーマットを、主旋律・対旋律をそれぞれ別行で
「音階0,音階1,音階2,音の長さ」の4値に変更（myPlayOrder==4は変更なし）。
SP\_Piano/SP\_Mokkin/SP\_flute側の\texttt{serialEvent()}を4値パースに対応させ、
受信した最大3音（0は休符）を同じdurationで同時発音するように修正した。

\subsection{Tick送信を8分音符基準から16分音符基準への変更}
\noindent
作業日時：2026年6月24日 \quad 作業時間：3時間
状態：完了（実機検証は未実施） \quad 進捗率：100\%（コード変更分）

\noindent
\textbf{依存関係・ブロッカー：}なし

\noindent
\textbf{作業内容：}

楽譜の解像度が8分音符単位から16分音符単位（104→208要素）に変更されたことに合わせて、
演奏システム全体の整合性を取るための修正を実施した。
\begin{itemize}
    \item \texttt{MA\_ptype1.ino}：Tick送信間隔の計算式\texttt{30000UL/BPM}（8分音符のms）を
    \texttt{15000UL/BPM}（16分音符のms）に変更（初期設定・BPM同期・通常送信の計3箇所＋最終待機の1箇所）。
    Tick終了判定を\texttt{tickCount>=104}から\texttt{tickCount>=208}に変更。
    \item \texttt{SA\_ptype2.ino}：myPlayOrder==2,3の演奏開始位置を新しいタイミング（tickCount=83で楽譜75要素目から開始、tickCount=139で138要素目に合流）に合わせて
    \texttt{startTick=37→82}, \texttt{indexOffset=32→74}、合流条件を\texttt{i2cReceiveCount>=66→>=139}に変更。
    ドラム専用処理（\texttt{handleDrum()}）はTickが2倍速くなった分、従来と同じ実時間での
    ハイハット／スネア交互発音間隔（4分音符ごと）を保つため判定を\texttt{\%2==0→\%4==0}に変更し、
    終了判定を\texttt{tickCount==103→207}に変更。
\end{itemize}
変更後、myPlayOrder==1/2,3・ドラムの境界Tick値（83, 138, 139, 207, 208等）について
index計算を手計算でトレースし、オーバーラン・ズレが無いことを机上で確認した。
実機への書き込み・Processing起動・通しでの演奏確認は未実施。

% ===== セクション3: 課題管理 =====
\section{課題管理}

\begin{table}[H]
    \centering
    \caption{課題一覧}
    \label{tab:issues}
    \begin{tabular}{lllllll}
        \toprule
        課題 & 発生原因 & 影響度 & 優先度 & 対策 & 期限 & 状態 \\
        \midrule
        楽器ごとの音量比設定 & 同一PC、複数Processingでの結合テスト & 中 & 中 & 1つの楽器につき1つのPCを使用する & 次回授業 & 未達成 \\
        \bottomrule
    \end{tabular}
\end{table}

\noindent
{\small ※影響度＝プロジェクトへのインパクト \quad ※優先度＝対応の緊急性}

% ===== セクション4: AI利用状況と検証 =====
\section{AI利用状況と検証（再現性重視）}

\subsection{利用内容}

\begin{itemize}
    \item \textbf{利用目的：}デバッグ（コンパイルエラー・演奏順不具合の原因特定）／実装（主旋律・対旋律の和音対応、楽譜の全面更新、ドラム機能追加、Tick解像度変更）／レビュー（Processing側ライブラリ・API不整合の修正、MP間の楽器順序整合性確認）
    \item \textbf{入力（プロンプト要旨）：}
    \begin{itemize}
        \item 修正前後のSA\_ptype1.inoの差分を提示し、問題点と原因の特定を依頼
        \item UNO R4 WiFi書き込み時のコンパイルエラーメッセージをそのまま提示し修正を依頼
        \item 「myPlayOrder==2,3を1番手の進行に65要素目で合流させたい」等、楽譜進行の仕様変更を自然言語で指示
        \item 「主旋律と対旋律を同時に演奏するようにSA・SPを変更せよ」「Drumslave\_ver2.inoを参考にmyPlayOrder==4のドラム機能を追加せよ」という要件付きの実装依頼
        \item Processing実行時のエラーメッセージ（ライブラリ未検出、メソッド不存在）をそのまま提示し修正を依頼
        \item 楽譜構造および送信データ形式変更に伴うシステム全体の整合性修正
    \end{itemize}
    \item \textbf{出力概要：}
    \begin{itemize}
        \item コンパイルエラーの原因（配列リファクタの取り残し、volatile String非対応）の説明と修正コード
        \item SA/SPのシリアル通信フォーマットを4値・3値に拡張する実装
        \item INST\_ID==4専用のドラム発音ロジック（handleDrum関数）
        \item processing.sound依存をddf.minimに置き換えたDrumProcessing\_new.pde
        \item Tick送信間隔・myPlayOrder開始・合流タイミング・ドラム発音間隔の修正コード
    \end{itemize}
\end{itemize}

\subsection{検証方法}

{\small ※第三者が再現できるレベルで記述}

\begin{itemize}
    \item \textbf{検証手法：}実機での回路作成・結合テスト
    \item \textbf{検証手順：}
    \begin{enumerate}
        \item 実機・実行環境でテストを行い、実際の動作やエラーコードを確認
    \end{enumerate}
\end{itemize}

\subsection{ファクトチェック結果}

\begin{itemize}
    \item \textbf{一致点：}
    \begin{itemize}
        \item 自分で設定した演奏順ごとの楽譜再生位置を基にAIが作成したコードでは、もりのくまさんの楽譜仕様通りの輪奏を行うことができた。
    \end{itemize}
    \item \textbf{不一致・疑義：}
    \begin{itemize}
        \item AIが作成した楽譜では、もりのくまさんとして演奏することができなかった。
    \end{itemize}
    \item \textbf{最終判断：}大枠はAIが作成したコードを採用し、適宜個人で修正を行なった。楽譜に関してはAIの出力は採用せず、全て個人で作成した。
    \item \textbf{根拠：}
    \begin{itemize}
        \item 機能や関数ごとにAIの出力を実機でのテストを行い、意図通りの動作をすること確認できた。AIの出力した楽譜では、音階や再生する音の数が異なり、意図通りの演奏ができなかったこと。
    \end{itemize}
\end{itemize}

% ===== セクション5: 次回計画 =====
\section{次回計画}

\begin{itemize}
    \item \textbf{目標：}
    \begin{itemize}
        \item 実機（Arduino UNO R4 WiFi×複数台、PC上のProcessing）でMA/MP/SA/SPを結合し、主旋律・対旋律・ドラムを含む全体演奏を通しで確認する
        \item 楽器Arduino一つにつき一台のPC（Processing）を使用したテストを行うことで、楽器ごとの適切な音量比を設定する。
    \end{itemize}
    \item \textbf{達成基準：}
    \begin{itemize}
        \item もりのくまさんの楽曲が冒頭から終端まで、ピアノ・木琴・フルート（主旋律＋対旋律）とドラムが破綻なく同時再生される
        \item 各楽器の適切な音量比を確定させる。
    \end{itemize}
    \item \textbf{期限：}次回授業
\end{itemize}

% ===== セクション6: その他 =====
\section{その他}

なし

% ===== セクション7: 付録 =====
\section{付録}

\textbf{GitHubリポジトリ}
https://github.com/Riarumogura/Hackathon1/tree/main/group7

\end{document}
